%=====================================================================
\chapter{Constraint Satisfaction}\label{ch:csp}
%=====================================================================
\locator{3}{Part III --- Constraints, Games and Plans}

\begin{outcomes}
By the end of this chapter you will be able to:
\begin{tight}
  \item \textbf{Formulate} a problem as variables, domains and constraints, and
        \textbf{say} when that formulation is better than a search over states.
  \item \textbf{Apply} backtracking search with the minimum-remaining-values and
        degree heuristics, and \textbf{explain} what each one orders.
  \item \textbf{Propagate} constraints by forward checking and by arc consistency
        (AC-3), and \textbf{measure} what each removes.
  \item \textbf{Explain} why ordering heuristics and propagation are a package,
        and why either alone can be worthless.
  \item \textbf{Recognise} when a constraint problem is over-constrained, and
        \textbf{report} which constraint is responsible.
\end{tight}
\end{outcomes}

\begin{prereq}
\chref{python} (states as values) and \chref{uninformed} (backtracking is
depth-first search with a commitment at each level). \chref{informed} is not
needed: this chapter gets its leverage from the \emph{structure} of the problem
rather than from a heuristic estimate of distance.
\end{prereq}

\begin{keyterms}
constraint satisfaction problem $\bullet$ variable $\bullet$ domain $\bullet$
constraint $\bullet$ constraint graph $\bullet$ assignment $\bullet$ partial
assignment $\bullet$ consistency $\bullet$ backtracking $\bullet$ minimum
remaining values $\bullet$ degree heuristic $\bullet$ least constraining value
$\bullet$ forward checking $\bullet$ arc consistency $\bullet$ AC-3 $\bullet$
over-constrained
\end{keyterms}

\section{A Sprint Plan as a Constraint Problem}

Twelve stories must be placed into four sprints. Each sprint has a capacity of
twelve story points. Some stories cannot start until others have finished. Some
pairs belong to the same developer and are too large to sit in one sprint
together. Find a placement that breaks nothing.

\chref{uninformed} could search this: a state is a partial placement, an action
adds one story. But that formulation throws away the thing that makes the problem
tractable. A general search algorithm sees an opaque state and an opaque goal test.
A \term{constraint satisfaction problem} exposes the structure --- which variable
each constraint touches --- and the algorithms in this chapter exploit exactly that.

\begin{definitionbox}[Constraint satisfaction problem]
A \term{CSP} is a triple:
\begin{tight}
  \item a set of \term{variables} $X_{1}, \ldots, X_{n}$;
  \item a \term{domain} $D_{i}$ of possible values for each variable;
  \item a set of \term{constraints}, each restricting the values that some subset
        of the variables may take simultaneously.
\end{tight}
A \term{solution} is a complete assignment that violates no constraint. A
constraint over two variables is \term{binary}; over more than two, $n$-ary.

The crucial property is that the goal test is not opaque: it is a conjunction of
constraints, each of which can be checked against a \emph{partial} assignment. That
is what allows a search to fail early instead of at the leaves.
\end{definitionbox}

\begin{examplebox}[The sprint problem, formulated]
\begin{tight}
  \item \emph{Variables:} the twelve stories \texttt{S1} \ldots \texttt{S12}.
  \item \emph{Domains:} $\{1,2,3,4\}$, the sprint each is placed in.
  \item \emph{Precedence} (binary): $\texttt{S1} < \texttt{S4}$,
        $\texttt{S2} < \texttt{S5}$, $\texttt{S3} < \texttt{S7}$,
        $\texttt{S6} < \texttt{S8}$, $\texttt{S9} < \texttt{S12}$,
        $\texttt{S10} < \texttt{S11}$, $\texttt{S4} < \texttt{S12}$.
  \item \emph{Same developer, too big together} (binary, ``not equal''):
        $\{\texttt{S1},\texttt{S7}\}$, $\{\texttt{S2},\texttt{S8}\}$,
        $\{\texttt{S4},\texttt{S6}\}$, $\{\texttt{S9},\texttt{S10}\}$,
        $\{\texttt{S7},\texttt{S12}\}$.
  \item \emph{Capacity} ($n$-ary): the points in each sprint total at most $12$.
        The twelve stories total $45$ points against a capacity of $48$, so there
        is almost no slack.
\end{tight}
There are $4^{12} = 16{,}777{,}216$ assignments and exactly $14$ solutions ---
about one in a million.
\end{examplebox}

\dsfig{%
\begin{tikzpicture}[font=\scriptsize,
  v/.style={circle, draw=#1, line width=0.9pt, fill=#1!10, text=#1!55!black,
            font=\tiny\bfseries, minimum size=0.78cm, inner sep=0pt},
  bef/.style={-{Stealth[length=1.8mm]}, primary!65, line width=0.9pt},
  apt/.style={accent!65, line width=0.9pt, dash pattern=on 3pt off 2pt}]

\node[v=primary]   (s1)  at (0,1.6)    {S1};
\node[v=secondary] (s2)  at (1.9,2.5)  {S2};
\node[v=primary]   (s3)  at (0,0)      {S3};
\node[v=greenhl]   (s4)  at (2.6,1.1)  {S4};
\node[v=secondary] (s5)  at (4.2,2.5)  {S5};
\node[v=greenhl]   (s6)  at (1.9,-1.1) {S6};
\node[v=primary]   (s7)  at (2.6,-0.2) {S7};
\node[v=secondary] (s8)  at (4.0,-1.1) {S8};
\node[v=goldhl]    (s9)  at (5.6,1.1)  {S9};
\node[v=goldhl]    (s10) at (6.6,-0.4) {S10};
\node[v=greenhl]   (s11) at (8.2,-0.4) {S11};
\node[v=primary]   (s12) at (5.0,0.3)  {S12};

\foreach \a/\b in {s1/s4,s2/s5,s3/s7,s6/s8,s9/s12,s10/s11,s4/s12}{
  \draw[bef] (\a) -- (\b);}
\foreach \a/\b in {s1/s7,s2/s8,s4/s6,s9/s10,s7/s12}{
  \draw[apt] (\a) -- (\b);}

\node[anchor=west, font=\scriptsize, text=primary] at (9.3,1.3)
  {\textbf{---$\blacktriangleright$} must finish earlier};
\node[anchor=west, font=\scriptsize, text=accent] at (9.3,0.85)
  {\textbf{- - -} cannot share a sprint};
\node[anchor=west, font=\scriptsize, text=neutral, text width=4.4cm]
  at (9.3,0.0)
  {Node colour is the owning developer. The \emph{capacity} constraint is
   $n$-ary and touches every variable at once, so it does not appear in this
   graph --- which matters in Section~\ref{sec:ac3}.};
\end{tikzpicture}}%
{The constraint graph. Twelve variables, twelve binary constraints, and one
$n$-ary capacity constraint that no edge can represent.}{csp-graph}

\section{Backtracking Search}

Assign one variable at a time; after each assignment check it against everything
already assigned; if nothing is consistent, undo the last choice and try the next
value. This is \term{backtracking search}, and it is depth-first search with the
goal test applied incrementally.

The incremental check is the whole gain. A general search must place all twelve
stories before it can test anything; backtracking discovers that \texttt{S4} cannot
go in sprint~1 the moment \texttt{S1} is placed there, and never builds any of the
$4^{10}$ assignments below that choice.

On the sprint problem, plain backtracking finds a solution after $123$ nodes and
$3{,}784$ consistency checks --- against $16{,}777{,}216$ complete assignments. The
structure was worth a factor of a hundred thousand before any heuristic was
applied.

\section{Which Variable Next, and Which Value}

Backtracking leaves two choices open, and they are where the remaining leverage is.

\textbf{Minimum remaining values (MRV).} Assign next the variable with the fewest
legal values left. The slogan is \emph{fail first}: if some variable is going to
run out of options, discovering it now costs one dead end, and discovering it ten
levels later costs a subtree.

\textbf{Degree heuristic.} When several variables tie on MRV, pick the one involved
in the most constraints on \emph{unassigned} variables. It is a tie-breaker, not a
strategy in its own right.

\textbf{Least constraining value.} Having chosen a variable, try first the value
that rules out fewest options for its neighbours. Note that this pulls the opposite
way from MRV --- fail first on variables, succeed first on values --- and both are
right, because you must try every variable but only need one value to work.

\begin{alertbox}[MRV on its own does nothing at all]
Here is a result from the lab that is worth more than the heuristics themselves.
Adding MRV to plain backtracking on this problem changes the cost by
\emph{nothing}: $3{,}784$ consistency checks and $123$ nodes, both with it and
without.

The reason is not that MRV is a bad idea. It is that MRV has no information to act
on. Every domain starts at four values and \emph{nothing ever shrinks them}, so
every unassigned variable ties forever and MRV falls back on the first variable in
the list --- which is what plain backtracking already did.

Domains shrink only when something propagates the consequences of an assignment.
So MRV and propagation are a package: neither is worth much without the other, and
together they are worth a factor of nine here. If you take one thing from this
chapter, take that.
\end{alertbox}

\section{Propagation: Forward Checking}

\term{Forward checking} is the smallest useful propagation. When a variable is
assigned, remove from every neighbour's domain the values now inconsistent with it.
If any domain becomes empty, the assignment has already failed --- backtrack at
once rather than discovering it later.

Two things happen at the same time. The obvious one is that dead ends are found
sooner. The important one is that domains now have different sizes, so MRV finally
has something to compare.

\dsfig{%
\begin{tikzpicture}[font=\scriptsize,
  dom/.style={rounded corners=2pt, draw=#1!60, fill=#1!8, text=#1!55!black,
              font=\tiny, minimum width=1.55cm, minimum height=0.46cm},
  hd/.style={font=\tiny\bfseries, text=neutral},
  a/.style={-{Stealth[length=1.8mm]}, neutral!70, line width=0.8pt}]

\node[hd] at (-1.35,0.75) {before};
\node[hd] at (-1.35,-0.95) {after};
\node[hd, text=primary] at (-1.35,0) {assign};

% before
\foreach \i/\n in {0/S4, 1/S6, 2/S7, 3/S12}{
  \node[dom=neutral] at (\i*1.85,0.75) {\n: 1 2 3 4};}
% the assignment
\node[rounded corners=2pt, fill=primary, text=white, font=\tiny\bfseries,
      minimum width=1.55cm, minimum height=0.46cm] at (0,0) {S1 := 1};
% after
\node[dom=greenhl] at (0,-0.95)    {S4: \phantom{1} 2 3 4};
\node[dom=neutral] at (1.85,-0.95) {S6: 1 2 3 4};
\node[dom=accent]  at (3.70,-0.95) {S7: \phantom{1} 2 3 4};
\node[dom=neutral] at (5.55,-0.95) {S12: 1 2 3 4};

\draw[a] (0,0.52) -- (0,0.18);
\draw[a] (0,-0.22) -- (0,-0.72);
\draw[a] (0.78,0) -- (3.70,-0.72);

\node[font=\tiny\itshape, text=greenhl!50!black, anchor=west] at (6.55,-0.72)
  {\texttt{S1}$<$\texttt{S4} removes sprint 1 from \texttt{S4}};
\node[font=\tiny\itshape, text=accent!55!black, anchor=west] at (6.55,-1.12)
  {\texttt{S1}, \texttt{S7} cannot share, so 1 goes from \texttt{S7}};

\node[rounded corners=3pt, fill=lightbg, draw=primary!30, line width=0.8pt,
      text=primary!70!black, font=\scriptsize, text width=13.2cm,
      align=flush left, inner sep=6pt] at (5.2,-2.35)
  {\textbf{And now MRV has something to say.} Before the assignment every domain
   had four values and MRV could not choose. After it, \texttt{S4} and \texttt{S7}
   have three and the others four --- so MRV picks a constrained variable, fails
   faster, and the two techniques begin to compound. Neither does this alone.};
\end{tikzpicture}}%
{Forward checking after one assignment. The pruning is the visible effect; giving
the ordering heuristic something to measure is the valuable one.}{csp-forward}

\section{Propagation: Arc Consistency}
\label{sec:ac3}

Forward checking only looks one step from the variable just assigned. \term{Arc
consistency} propagates further.

\begin{definitionbox}[Arc consistency and AC-3]
An arc $X_{i} \to X_{j}$ is \term{consistent} if for \emph{every} value in
$D_{i}$ there is \emph{some} value in $D_{j}$ satisfying the binary constraint
between them. A value in $D_{i}$ with no partner can never appear in a solution
and is removed.

\term{AC-3} enforces this everywhere: keep a queue of arcs; take one, prune $D_{i}$
against $D_{j}$; if anything was removed, re-queue every arc \emph{into} $X_{i}$,
because those may now have lost partners. Stop when the queue empties, or fail if
a domain empties.
\end{definitionbox}

Run as a preprocessing step on the sprint problem, before any search, AC-3 narrows
every one of the twelve domains: stories at the head of a precedence chain lose the
last sprint, stories at the tail lose the first. That is real information obtained
without a single assignment.

\begin{alertbox}[What AC-3 cannot see, and it is the important constraint]
Look again at \dsref{csp-graph}: the capacity constraint has no edge, because it is
$n$-ary and a constraint graph draws binary constraints. AC-3 works on arcs, so it
never considers capacity at all.

This matters, because capacity is the constraint that actually makes this problem
hard --- $45$ points into $48$ of capacity. Arc consistency prunes the precedence
and same-developer structure beautifully and leaves the binding constraint
untouched. The measurement in Table~\ref{tab:csp} shows the consequence: AC-3
preprocessing costs $224$ extra consistency checks and removes \emph{no} search
nodes at all.

The general lesson is not that AC-3 is bad. It is that \emph{propagation only
propagates what you have expressed in a form it can read}. If the binding
constraint is $n$-ary, you need a generalised consistency algorithm, or you must
re-express the constraint --- and knowing which situation you are in is the
engineering judgement here.
\end{alertbox}

\section{The Measurement}

All five configurations solve the same instance and find a solution. What differs
is how much work it costs.

\smallskip
\noindent\begin{longtable}{@{}L{6.0cm}L{2.3cm}L{2.0cm}L{4.4cm}@{}}
\toprule
\rowcolor{primary!15}
\textbf{Configuration} & \textbf{Checks} & \textbf{Nodes} & \textbf{Comment} \\
\midrule
\endhead
Plain backtracking & 3{,}784 & 123 & The baseline. Already $10^{5}$ better than
enumeration \\
\quad + MRV & 3{,}784 & 123 & \textbf{No change.} Nothing shrinks the domains, so
MRV never has a preference \\
\quad + MRV + degree & 4{,}333 & 436 & \textbf{Worse.} The tie-break imposes an
ordering that happens to be unlucky here \\
\quad + MRV + degree + forward checking & \textbf{415} & \textbf{13} &
\textbf{9$\times$ fewer checks.} Propagation gives MRV information, and the two
compound \\
\quad + \ldots + AC-3 preprocessing & 639 & 13 & Same nodes, more checks: AC-3
cannot see the capacity constraint \\
\bottomrule
\caption{Twelve stories into four sprints, counting every constraint check. A
solution exists (there are $14$) and every configuration finds one.}
\label{tab:csp}
\end{longtable}

Three conclusions, and the second is the one students get wrong in examinations.

\textbf{Structure beats search.} Before any heuristic, formulating the problem as a
CSP turned $16.8$ million assignments into $123$ nodes.

\textbf{An ordering heuristic can make things worse.} The degree tie-breaker
tripled the node count. Heuristics are informed guesses about a specific problem,
not improvements in general, and the only way to know is to measure on your
instance.

\textbf{Propagation is what pays.} Forward checking cut the work ninefold --- and
about half of that was the pruning and half was making MRV functional.

\section{When There Is No Solution}

An over-constrained problem is the common case in practice: the backlog does not
fit. Backtracking then explores the whole space and reports failure, which is true
and useless. A delivery manager needs to know \emph{which} constraint to relax.

Two cheap and honest answers. Record the constraint that emptied a domain most
often during propagation --- that is usually the binding one. Or re-run with each
constraint removed in turn and report which removals make the problem solvable.
For this instance, raising the capacity from $12$ to $13$ turns $14$ solutions into
hundreds, which tells you at once that capacity, not precedence, is what binds.

\begin{worked}[Tracing forward checking through the first three assignments]
Start with every domain $\{1,2,3,4\}$. Take the variables in order.

\smallskip
\textbf{Assign \texttt{S1} $:=$ 1.} Forward checking examines \texttt{S1}'s
neighbours in \dsref{csp-graph}:
\begin{tight}
  \item $\texttt{S1} < \texttt{S4}$: \texttt{S4} must be later than $1$, so
        $D_{\texttt{S4}} = \{2,3,4\}$.
  \item $\texttt{S1} \ne \texttt{S7}$: $D_{\texttt{S7}} = \{2,3,4\}$.
\end{tight}
Capacity: sprint~1 now holds $3$ points of $12$. Fine.

\smallskip
\textbf{MRV now chooses.} \texttt{S4} and \texttt{S7} have three values; everything
else has four. It picks one of those two --- and note that this is the first moment
in the entire search at which MRV has expressed a preference. Say \texttt{S4}.

\smallskip
\textbf{Assign \texttt{S4} $:=$ 2.} Propagate:
\begin{tight}
  \item $\texttt{S4} < \texttt{S12}$: $D_{\texttt{S12}} = \{3,4\}$.
  \item $\texttt{S4} \ne \texttt{S6}$: $D_{\texttt{S6}} = \{1,3,4\}$.
\end{tight}
Capacity: sprint~2 holds $8$ of $12$; only $4$ points remain there, which will
matter shortly.

\smallskip
\textbf{MRV chooses \texttt{S12}}, now down to two values --- the most constrained
variable in the problem, and it was reached in two steps rather than by luck.

\smallskip
\textbf{Assign \texttt{S12} $:=$ 3.} Propagate:
$\texttt{S9} < \texttt{S12}$ gives $D_{\texttt{S9}} = \{1,2\}$;
$\texttt{S7} \ne \texttt{S12}$ gives $D_{\texttt{S7}} = \{2,4\}$.

\smallskip
\textbf{What has been bought.} Three assignments, and four other variables have
already lost values. Plain backtracking at this depth has pruned nothing and will
discover the same facts one dead end at a time, hundreds of nodes later. The lab
counts the difference: $13$ nodes against $123$.
\end{worked}

\begin{pitfall}
\begin{tight}
  \item \textbf{Expecting MRV to help on its own.} Without propagation the domains
        never shrink, so MRV has nothing to order by. Measured here: exactly zero
        improvement.
  \item \textbf{Assuming a heuristic cannot hurt.} The degree tie-breaker tripled
        the node count on this instance.
  \item \textbf{Expecting AC-3 to enforce an $n$-ary constraint.} It works on arcs.
        Capacity is invisible to it, and capacity is what binds here.
  \item \textbf{Re-running AC-3 from scratch at every node.} Propagate
        incrementally from the variable just assigned; a full re-run usually costs
        more than it saves, as the last row of Table~\ref{tab:csp} shows even for
        a single preprocessing pass.
  \item \textbf{Reporting ``no solution'' and stopping.} True and useless. Report
        which constraint is binding.
  \item \textbf{Counting nodes and calling it performance.} A node with
        propagation costs more than a node without. Count consistency checks, or
        time, or both --- the same mistake as \chref{local}'s evaluation budget.
\end{tight}
\end{pitfall}

%---------------------------------------------------------------------
\begin{lab}[Sprint Scheduling as a CSP]
The instance of \dsref{csp-graph}, solved five ways, counting every constraint
check. Reproduces Table~\ref{tab:csp}.
\begin{lstlisting}[language=Python]
"""Chapter 7 lab -- sprint scheduling as a constraint problem.

Twelve stories into four sprints under precedence, same-developer and
capacity constraints. The headline result is the second row: adding
MRV to plain backtracking changes nothing at all, because nothing has
shrunk the domains for it to compare.
"""

import itertools

# story -> (points, owner)
STORIES = {
    "S1": (3, "Ayesha"),  "S2": (5, "Bilal"), "S3": (2, "Ayesha"),
    "S4": (8, "Chen"),    "S5": (3, "Bilal"), "S6": (2, "Chen"),
    "S7": (5, "Ayesha"),  "S8": (3, "Bilal"), "S9": (5, "Dina"),
    "S10": (3, "Dina"),   "S11": (2, "Chen"), "S12": (4, "Ayesha"),
}
NAMES = list(STORIES)
SLOTS = (1, 2, 3, 4)
CAP = 12

# binary: a must finish in an earlier sprint than b
BEFORE = [("S1", "S4"), ("S2", "S5"), ("S3", "S7"), ("S6", "S8"),
          ("S9", "S12"), ("S10", "S11"), ("S4", "S12")]
# binary: same developer, too much combined work for one sprint
APART = [("S1", "S7"), ("S2", "S8"), ("S4", "S6"), ("S9", "S10"),
         ("S7", "S12")]


class Counter:
    def __init__(self):
        self.checks = 0
        self.nodes = 0


def binary_ok(a, va, b, vb, c):
    """Is the pair (a=va, b=vb) consistent? One consistency check."""
    c.checks += 1
    if (a, b) in BEFORE and not va < vb:
        return False
    if (b, a) in BEFORE and not vb < va:
        return False
    if ((a, b) in APART or (b, a) in APART) and va == vb:
        return False
    return True


def capacity_ok(assign, c):
    """The n-ary constraint. No arc can express this."""
    c.checks += 1
    load = dict.fromkeys(SLOTS, 0)
    for n, v in assign.items():
        load[v] += STORIES[n][0]
    return all(load[s] <= CAP for s in SLOTS)


def consistent(assign, var, val, c):
    for other, oval in assign.items():
        if not binary_ok(var, val, other, oval, c):
            return False
    trial = dict(assign)
    trial[var] = val
    return capacity_ok(trial, c)


def select_var(assign, domains, mrv, degree):
    un = [v for v in NAMES if v not in assign]
    if not mrv:
        return un[0]
    fewest = min(len(domains[v]) for v in un)
    tied = [v for v in un if len(domains[v]) == fewest]
    if len(tied) == 1 or not degree:
        return tied[0]

    def deg(v):
        return sum(1 for (a, b) in BEFORE + APART
                   if (a == v and b not in assign)
                   or (b == v and a not in assign))

    return max(tied, key=deg)


def forward_check(var, val, domains, c):
    """Prune the neighbours. None means a domain emptied."""
    nd = {k: list(v) for k, v in domains.items()}
    nd[var] = [val]
    for other in NAMES:
        if other == var:
            continue
        keep = [x for x in nd[other] if binary_ok(var, val, other, x, c)]
        if not keep:
            return None
        nd[other] = keep
    return nd


def ac3(domains, c):
    """Arc consistency over the BINARY constraints only."""
    nd = {k: list(v) for k, v in domains.items()}
    arcs = []
    for (a, b) in BEFORE + APART:
        arcs += [(a, b), (b, a)]
    queue = list(arcs)
    while queue:
        xi, xj = queue.pop(0)
        keep = [vi for vi in nd[xi]
                if any(binary_ok(xi, vi, xj, vj, c) for vj in nd[xj])]
        if len(keep) != len(nd[xi]):
            nd[xi] = keep
            if not nd[xi]:
                return None
            queue += [(a, b) for (a, b) in arcs if b == xi]
    return nd


def solve(mrv=False, degree=False, fc=False, use_ac3=False):
    c = Counter()
    domains = {v: list(SLOTS) for v in NAMES}
    if use_ac3:
        domains = ac3(domains, c)
        if domains is None:
            return None, c
    found = []

    def bt(assign, domains):
        c.nodes += 1
        if len(assign) == len(NAMES):
            found.append(dict(assign))
            return True
        var = select_var(assign, domains, mrv, degree)
        for val in domains[var]:
            if consistent(assign, var, val, c):
                assign[var] = val
                if fc:
                    nd = forward_check(var, val, domains, c)
                    if nd is not None and bt(assign, nd):
                        return True
                elif bt(assign, domains):
                    return True
                del assign[var]
        return False

    bt({}, domains)
    return (found[0] if found else None), c


def count_solutions():
    """Brute force over all 4^12 assignments, for the record."""
    c, n = Counter(), 0
    for combo in itertools.product(SLOTS, repeat=len(NAMES)):
        a = dict(zip(NAMES, combo))
        if all(a[x] < a[y] for x, y in BEFORE) \
           and all(a[x] != a[y] for x, y in APART) \
           and capacity_ok(a, c):
            n += 1
    return n


def show(sol):
    by = {s: [] for s in SLOTS}
    for n, v in sol.items():
        by[v].append(n)
    out = []
    for s in SLOTS:
        pts = sum(STORIES[n][0] for n in by[s])
        out.append("sprint %d: %-18s %2d pts"
                   % (s, ",".join(sorted(by[s], key=lambda x: int(x[1:]))),
                      pts))
    return "\n".join(out)


if __name__ == "__main__":
    total = sum(p for p, _ in STORIES.values())
    print("stories %d   sprints %d   assignments %d"
          % (len(NAMES), len(SLOTS), len(SLOTS) ** len(NAMES)))
    print("points %d against capacity %d -- almost no slack"
          % (total, CAP * len(SLOTS)))
    print("solutions:", count_solutions())
    print()
    print("configuration                    checks  nodes")
    print("-" * 50)
    runs = [
        ("plain backtracking           ", dict()),
        ("+ MRV                        ", dict(mrv=True)),
        ("+ MRV + degree               ", dict(mrv=True, degree=True)),
        ("+ MRV + degree + FC          ", dict(mrv=True, degree=True,
                                                fc=True)),
        ("+ MRV + degree + FC + AC-3   ", dict(mrv=True, degree=True,
                                                fc=True, use_ac3=True)),
    ]
    res = {}
    for name, kw in runs:
        sol, c = solve(**kw)
        assert sol is not None, name
        res[name.strip()] = (c.checks, c.nodes)
        print("%s %7d %6d" % (name, c.checks, c.nodes))

    plain = res["plain backtracking"]
    mrv_only = res["+ MRV"]
    fc = res["+ MRV + degree + FC"]
    ac = res["+ MRV + degree + FC + AC-3"]

    # 1. MRV alone changes NOTHING: no propagation, no information
    assert mrv_only == plain
    print()
    print("MRV alone: %d checks, exactly as without it. Nothing has"
          % mrv_only[0])
    print("shrunk a domain, so every variable ties and MRV cannot")
    print("express a preference.")

    # 2. forward checking is the win, and it is large
    assert fc[0] * 5 < plain[0]
    print()
    print("forward checking: %d -> %d checks (%.0fx), %d -> %d nodes"
          % (plain[0], fc[0], plain[0] / fc[0], plain[1], fc[1]))

    # 3. AC-3 cannot see the capacity constraint, so it buys no nodes
    assert ac[1] == fc[1] and ac[0] > fc[0]
    print()
    print("AC-3 preprocessing: %d extra checks, %d fewer nodes."
          % (ac[0] - fc[0], fc[1] - ac[1]))
    print("It prunes the binary structure and never sees capacity,")
    print("which is the constraint that actually binds.")

    # what AC-3 removes before any search
    c = Counter()
    d0 = {v: list(SLOTS) for v in NAMES}
    d1 = ac3(d0, c)
    narrowed = [v for v in NAMES if d1[v] != d0[v]]
    print()
    print("AC-3 narrowed %d of %d domains before any assignment:"
          % (len(narrowed), len(NAMES)))
    for v in narrowed[:4]:
        print("   %-4s %s -> %s" % (v, d0[v], d1[v]))
    print("   ...")

    print()
    sol, _ = solve(mrv=True, degree=True, fc=True)
    print(show(sol))
\end{lstlisting}
Two experiments afterwards. Raise \texttt{CAP} to $13$ and count the solutions
again --- the problem stops being interesting, which is what ``the binding
constraint'' means. Then lower it to $11$ and confirm that the search reports
failure after exploring the whole space; that is the over-constrained case of the
last section, and it is the one a delivery manager will actually bring you.
\end{lab}

%---------------------------------------------------------------------
\begin{checkpoint}
\begin{tightnum}
  \item Adding MRV to plain backtracking on this problem changed the number of
        consistency checks by exactly zero. Explain why, and state the one change
        that makes MRV worth having.
  \item AC-3 narrowed every domain in the problem and removed no search nodes at
        all. Explain how both of those can be true at once.
  \item A delivery manager reports that the backlog does not fit and asks what to
        do. Your solver returns ``no solution''. Describe two ways to turn that
        into a useful answer, and say which constraint you would expect to be
        binding here.
\end{tightnum}
\end{checkpoint}

%---------------------------------------------------------------------
\begin{chaptersummary}
\begin{tight}
  \item A \term{CSP} is variables, domains and constraints. Its advantage over a
        general search is that the goal test is not opaque: constraints can be
        checked against a \emph{partial} assignment, so the search fails early.
        That alone turned $16.8$ million assignments into $123$ nodes.
  \item \term{Backtracking search} assigns one variable at a time and undoes the
        last choice on failure.
  \item \term{MRV} assigns the most constrained variable next --- fail first. The
        \term{degree heuristic} breaks ties; \term{least constraining value}
        orders values the other way, and both are right.
  \item \term{Forward checking} prunes neighbours' domains after each assignment.
        It cut the work ninefold here --- half by pruning, half by giving MRV
        something to measure. \emph{Ordering and propagation are a package.}
  \item \term{Arc consistency} (AC-3) propagates further, over binary constraints
        only. On this problem it narrowed all twelve domains and removed no nodes,
        because the binding constraint is $n$-ary and invisible to it.
  \item A heuristic is a guess about a particular problem, not a general
        improvement: the degree tie-breaker tripled the node count. Measure on
        your instance, counting consistency checks rather than nodes.
  \item For an over-constrained problem, report which constraint binds. ``No
        solution'' is true and useless.
\end{tight}
\end{chaptersummary}

\begin{reviewq}
  \item \textbf{(MCQ)} The minimum-remaining-values heuristic chooses:
        (a)~the value ruling out fewest options (b)~the variable with the fewest
        legal values left (c)~the variable in the most constraints (d)~the
        cheapest variable.
  \item \textbf{(MCQ)} Forward checking prunes: (a)~all variables
        (b)~the neighbours of the variable just assigned (c)~only the variable
        just assigned (d)~nothing until a domain empties.
  \item \textbf{(MCQ)} AC-3 operates on: (a)~$n$-ary constraints (b)~binary
        constraints (c)~the objective function (d)~the search tree.
  \item \textbf{(MCQ)} An arc $X_{i} \to X_{j}$ is consistent when:
        (a)~every value of $X_{i}$ has some consistent partner in $X_{j}$
        (b)~some value of $X_{i}$ has a partner (c)~$D_{i} = D_{j}$
        (d)~$X_{i}$ is assigned.
  \item \textbf{(Short)} Why does MRV need propagation to be useful? Use the
        measured result from this chapter in your answer.
  \item \textbf{(Short)} Explain why MRV picks the most constrained variable while
        least-constraining-value picks the least constraining value, and why that
        is not a contradiction.
  \item \textbf{(Short)} Give an $n$-ary constraint from a project you know, and
        say what goes wrong if you expect arc consistency to enforce it.
  \item \textbf{(Applied)} Take the constraint graph of \dsref{csp-graph}. By hand,
        assign $\texttt{S10} := 3$ and apply forward checking. Which domains
        change, and to what? Then say which variable MRV would select next, and
        whether the degree heuristic would be consulted.
\end{reviewq}
